% Copyright (C) 2003 Johan Hoffman and Anders Logg.
% Licensed under the GNU GPL Version 2.

\documentclass[12pt]{article}

%--- Packages ---
\usepackage[latin1]{inputenc}
\usepackage{graphicx}
\usepackage{fancyhdr}

\setlength{\parindent}{0pt}
\setlength{\parskip}{12pt}

%--- Header an footer ---
\lhead{\sf DOLFIN}
\rhead{\sf User Manual}
\pagestyle{fancy}
\renewcommand{\footrulewidth}{1pt}
\renewcommand{\headrulewidth}{1pt}

\begin{document}

%--- Title page ---

\thispagestyle{empty}

\Large
\textbf{DOLFIN User Manual} \\
\vspace{-0.5cm}
\hrule
\hfill
\small
\today

\vspace{3cm}

\begin{center}
\includegraphics[width=6cm]{eps/dolfin.eps}
\end{center}

\vfill

\normalsize
\textbf{Hoffman, Logg, et al.}
\hrule

\newpage

%--- List of authors ---

\ 
\vfill
This manual has been written by: \\
Johan Hoffman, Johan Jansson, and Anders Logg.
\thispagestyle{empty}

\cleardoublepage

%--- Table of contents ---

\tableofcontents
\newpage

%--- New chapter ---
\cleardoublepage
\section{Introduction}

This is a first draft for a DOLFIN manual. Contributions are most welcome.

%--- New chapter ---
\cleardoublepage
\section{Installation}

In preparation.

%--- New chapter ---
\cleardoublepage
\section{Linear algebra}

In preparation.

%--- New chapter ---
\cleardoublepage
\section{Grid management}

In preparation.

%--- New chapter ---
\cleardoublepage
\section{The log system}

The purpose of the log system is to provide a simple and
clean interface for logging messages, including warnings
and errors.

The following functions / macros are provided for logging:

\begin{verbatim}
  dolfin_info();
  dolfin_debug();
  dolfin_warning();
  dolfin_error();
  dolfin_assert();
\end{verbatim}

Examples of usage:

\begin{verbatim}
  dolfin_info("Created vector of size %d.", x.size());
  dolfin_debug("Opened file");
  dolfin_warning("Unknown cell type.");
  dolfin_error("Out of memory.");
  dolfin_assert(i < m);
\end{verbatim}

Note that in order to pass additional arguments to the last
three functions (which are really macros, in order to
automatically print information about file names, line numbers
and function names), the variations \texttt{dolfin\_debug1()},
\texttt{dolfin\_debug2()} and so on, must be used.

As an alternative to \texttt{dolfin\_info()}, C++ style output to cout
(\texttt{dolfin::cout}, and not \texttt{std::cout}) can be used. These messages
will be delivered to the same destination as messages by use
of the function \texttt{dolfin\_info()}.

Examples of usage:

\begin{verbatim}
  cout << "Assembling matrix: " << A << endl;
  cout << "Refining grid: " << grid << endl;
\end{verbatim}

The \texttt{dolfin\_assert()} macro should be used for simple tests that
may occur often, such as checking indexing in vectors. The check
is turned on only if DEBUG is defined.

To notify progress by a progress session, use the class
Progress.

Examples of usage:

\begin{verbatim}
  Progress p("Assembling", grid.noCells());
  
  for (CellIterator c(grid); !c.end(); ++c) {
    ...
    p++;
  }
\end{verbatim}

Progress also supports the following usage:

\begin{verbatim}
  p = i;    // Specify step number
  p = 0.5;  // Specify percentage
  p.update(t/T, "Time is t = %f", t);
\end{verbatim}

%--- New chapter ---
\cleardoublepage
\section{Handling parameters}

In preparation.

%--- New chapter ---
\cleardoublepage
\section{Writing a new module / solver}

In preparation.

%--- New chapter ---
\cleardoublepage
\section{Installing DOLFIN}

In preparation.

%--- New chapter ---
\cleardoublepage
\section{Contributing to DOLFIN}

If you have created a new module, fixed a bug somewhere, or have done
a small change which you want to contribute to DOLFIN, then
the best way to do so is to send a patch to the maintainer. A patch
is a file which describes how to transform a file or directory structure
into another. The patch is built by comparing a version which both parties
have against the modified version which only you have.

The tool used to create a patch is called \texttt{diff} and the tool
used to apply the patch is called \texttt{patch}. These tools are free
software and are standard in most UNIX-style operating systems.

\subsection{Creating a patch using CVS}

The simplest way to create a patch is through CVS. Assuming that you
have previously checked out a CVS version of DOLFIN and that you are positioned at
the top of the DOLFIN source tree, a patch can be created using the command

\texttt{\$ cvs diff > dolfin-logg-2004-01-16.patch}

Choose a suitable name for the patch to simplify identification, for example your name and
today's date.

\subsection{Creating a patch without using CVS}

Example:

Let's say you have started from the DOLFIN release 0.3.10. You have a
directory structure under \texttt{dolfin-0.3.10} where you have made
some modifications to some files which you think could be useful to
other DOLFIN users.

\begin{enumerate}
\item{

Clean up your modified directory structure to remove temporary and binary
files which will be rebuilt anyway:

\texttt{\$ make clean}
}
\item{

Rename the directory to something else:

\texttt{\$ mv dolfin-0.3.10 dolfin-0.3.10mod}
}
\item{

Unpack the version of DOLFIN which you started from:

\texttt{\$ tar xzvf dolfin-0.3.10.tgz}
}

\item{

You should now have two DOLFIN directory structures in your current directory:

\begin{verbatim}
$ ls
dolfin-0.3.10
dolfin-0.3.10mod
\end{verbatim}
}

\item{

Use the \texttt{diff} tool to create the patch:

\begin{verbatim}
$ diff -u --new-file --recursive dolfin-0.3.10 dolfin-0.3.10mod >
mypatch.patch
\end{verbatim}

``\texttt{-u}'' means ``unified'', a format in which to describe the patch.

``\texttt{--new-file}'' accepts the creation of new files in addition to modifications.

``\texttt{--recursive}'' means to recursively compare files through the
directory structure.


}
\item{

The patch now exists as \texttt{mypatch.patch} and can be distributed
to other people who have dolfin-0.3.10 to easily create your modified
version.  If the patch is large, compressing it with for example
\texttt{gzip} is advisable.

}

\end{enumerate}

%A patch can also be created relative to a CVS version of DOLFIN. Make
%sure to notify receivers of the patch exactly which CVS version it is
%however (date and time of checkout should suffice), since the patch
%should be applied to the same version used to build it. Patching is
%possible against a different source, but conflicts can arise, and even
%though there are heuristics which take care of most simple cases, the
%heuristics can fail, and there can be conflicts which are
%theoretically impossible to resolve.

\subsubsection{Applying a patch}

Example:

Let's say the patch has been built relative to DOLFIN release 0.3.10.

\begin{enumerate}
\item{

Unpack the version of DOLFIN which the patch is built relative to:

\texttt{\$ tar xzvf dolfin-0.3.10.tgz}
}

\item{

Check that you have the patch \texttt{mypatch.patch} and the DOLFIN
directory structure in the current directory.

\begin{verbatim}
$ ls
dolfin-0.3.10
mypatch.patch
\end{verbatim}
}

\item{

Enter the DOLFIN directory structure:

\begin{verbatim}
$ cd dolfin-0.3.10
\end{verbatim}
}

\item{

Apply the patch:

\begin{verbatim}
$ patch -p1 < ../mypatch.patch
\end{verbatim}
}

``\texttt{-p1}'' strips the leading directory from the filename
references in the patch, to match the fact that we are applying the
patch from inside the directory.

\item{

The modified version now exists as \texttt{dolfin-0.3.10}

}

\end{enumerate}

\end{document}
